\documentclass[10pt,a4paper]{amsart}
\usepackage[margin=19mm]{geometry}
\usepackage{amsmath,amssymb,mathtools}
\usepackage[scr=rsfs,cal=euler]{mathalfa}
\usepackage{xfrac}
\usepackage[unicode,hidelinks,bookmarks=false]{hyperref}
\newcommand{\ie}{i.e.\ }
\newcommand{\cf}{cf.\ }
\newcommand{\resp}{respectively\ }
\newcommand{\Subsection}{\subsection}
\providecommand{\nobreakdash}{\nobreak\hbox{-}}
\setlength{\emergencystretch}{2em}
\multlinegap0pt
\usepackage{bm}
\usepackage{bbm}
\usepackage{dsfont}
\usepackage{xspace}
\usepackage{cancel}
\usepackage{mathtools}
\usepackage{amsthm}
\makeatletter
\@ifundefined{lemma}{\newtheorem{lemma}{Lemma}}{}
\makeatother
\title[Non-compact Ricci Solitons of Finite Volume with Potential F]{Non-compact Ricci Solitons of Finite Volume with Potential Field of Constant Length}
\author{Hemangi Madhusudan Shah, Sharief Deshmukh,, Mohammad Aqib}
\date{Source: arXiv:2607.22394v1 (CC BY 4.0)}
\begin{document}
\maketitle

\noindent\emph{Source and licence.} Non-compact Ricci Solitons of Finite Volume with Potential Field of Constant Length. Version arXiv:2607.22394v1 (CC BY 4.0), distributed under the Creative Commons Attribution 4.0 International licence, as recorded in the arXiv licence metadata for this exact version and shown on the abstract page https://arxiv.org/abs/2607.22394v1. Full rights notice: licenses/arxiv-2607.22394-CC-BY-4.0.txt.\par\smallskip

\noindent\textit{Editorial scope.} Counted unit: the Lemma whose source label is \texttt{alpha-cosym} in the article (source0.tex lines 646--670, source label \texttt{alpha-cosym}), delivered together with the article's own complete original proof of it (source0.tex lines 672--753, one \texttt{proof} environment of 82 lines). No mathematics is written, repaired or re-derived: statement and proof are the article's own text line for line. Required context retained uncounted: e1 (lines 108--111); e2 (lines 123--125); e5 (lines 167--173); e9 (lines 326--328). Every definition, display and label of those lines is retained unchanged. Omitted from this package and not redistributed: 1--107, 112--122, 126--166, 174--325, 329--645, 671--671, 754--1274. Nothing inside a retained excerpt is omitted or altered: every retained body is delivered exactly as the article's lines are written, so no omission receipt of any kind accompanies this package. Citations: the retained excerpts carry no citation command at all, so no bibliography entry of the article is transcribed and no reference is redistributed. Reference adaptation: none was needed and none was made; every reference inside the delivered unit resolves inside it, so no unresolved reference or citation is produced. Printed slips kept literal: no printed word, symbol, hypothesis or number of the counted statement or of its proof was altered. Layout and class: the article's own class is \texttt{article} with packages including amsmath, amsthm, amsfonts, latexsym, amscd, amssymb, enumerate, hyperref, xcolor, soul, xparse; the wrapper is the lane's pinned \texttt{amsart} editorial wrapper at 10pt on A4, which restates the theorem-like environments the retained text needs and adds the article's own macro definitions that the retained text needs (none), with extra wrapper packages (bm, bbm, dsfont, xspace, cancel, mathtools). The delivered numbering is the unit's own. Assets: no retained span contains a figure environment, a graphics-inclusion command, a PDF-page inclusion, a table or a TikZ picture; no image file is copied into this package, the package declares no asset dependency and no image file is redistributed with it. Licence: version arXiv:2607.22394v1 of this article is distributed under the Creative Commons Attribution 4.0 International licence, as recorded in the arXiv licence metadata for this exact version and shown on the abstract page https://arxiv.org/abs/2607.22394v1. A full rights notice is delivered at licenses/arxiv-2607.22394-CC-BY-4.0.txt. Characters: every retained excerpt is pure ASCII, so the delivered engine, XeLaTeX with its default Latin Modern text font, reports no missing character.
\begin{equation}\label{e1}
{\nabla} _{U}X=\lambda U-QU+\phi U\text{,\quad }U\in \Gamma \left(
TM^{n}\right) \text{,}  
\end{equation}

\begin{equation}\label{e2}
\lambda X-QX -\phi X=0\text{.}  
\end{equation}

\begin{eqnarray}\label{e5}
\left\Vert Q\right\Vert ^{2}-\frac{1}{n}S^{2} &=&-\frac{1}{n}\left(
S-n\lambda \right) ^{2}+\frac{1}{2}X(S)   \\
&&-\sum\limits_{i}g\left( X,\left( \nabla \phi \right) \left(
u_{i},u_{i}\right) \right) -\left\Vert \phi \right\Vert ^{2}\text{.} 
\nonumber
\end{eqnarray}

\begin{equation}\label{e9}
\left\Vert Q\right\Vert ^{2}\geq \frac{1}{n}S^{2}  
\end{equation}

\begin{lemma}\label{alpha-cosym}
Let $(M^n,g,\xi,\lambda)$ be an almost $\alpha$-cosymplectic
Ricci soliton with potential vector field $\xi$, the Reeb vector
field of the almost $\alpha$-cosymplectic manifold
$(M,g,\varphi,\xi,\eta)$ and scalar curvature $S$.
Then the soliton function satisfies
\[
n\lambda^2-(1+2S)\lambda+\|Q\|^2=0.
\]
Consequently,
\begin{equation}\label{lambdaeq}
\lambda=
\frac{(1+2S)\pm\sqrt{(1+2S)^2-4n\|Q\|^2}}{2n}.
\end{equation}
Moreover,
\[
1+2S>0
\qquad\text{and}\qquad
\lambda\ge0.
\]
Hence, the soliton is non-expanding (that is, shrinking or steady).
Furthermore, if $\lambda=0$, then $\xi$ is parallel and therefore
$M^3=\mathbb R^3/\Gamma_1$, while in general
$M^n=N^{n-1}\times(a,b)$, where $N$ has finite volume.
\end{lemma}

\begin{proof}
Note that $\eta$ is closed and therefore (\ref{e1})
and (\ref{e2}) imply that
\begin{eqnarray}\label{eq:nablaxi}
\nabla_{X} \xi = \lambda X - Q X \\ \nonumber
Q(\xi) = \lambda \xi.
\end{eqnarray}
Using the above equation, we compute
\begin{eqnarray}\label{curv}
R(U,V)\xi = - (\nabla Q)(U,V) + (\nabla Q)(V,U),
\end{eqnarray}
which, on contraction, yields
\begin{eqnarray}\label{e11}
\operatorname{Ric}(V,\xi)  = -\frac{1}{2} V(S) + V(S) =  \frac{1}{2} V(S)
\end{eqnarray}
that is,
\begin{eqnarray}\label{Qxi}
Q(\xi) = \frac{1}{2} \nabla S.
\end{eqnarray}
Note that \eqref{e5} together with the fact that
$\xi$ is closed, takes the form
\begin{eqnarray}\label{m1}
||Q||^2 - \frac{S^2}{n}  =
- \frac{1}{n} \left(S - n \lambda \right)^2 + \frac{1}{2}
\xi(S).
\end{eqnarray}
Now \eqref{eq:nablaxi} and \eqref{e11} respectively imply
$$\operatorname{Ric}(\xi, \xi) = \lambda \;\mbox{and} \;
\operatorname{Ric}(\xi, \xi) = \frac{1}{2} \xi(S).$$
Thus,
\begin{eqnarray}\label{m2}
\frac{1}{2} \xi(S) = \lambda.
\end{eqnarray}
Plugging (\ref{m2}) into (\ref{m1}) yields
\[
n\lambda^{2}-(1+2S)\lambda+\Vert Q\Vert^{2}=0 ,
\]
whose roots are given by \eqref{lambdaeq}.

It remains to show that $\lambda\geq0$. Since the quadratic equation admits the real root $\lambda$,
its discriminant must be non-negative at every point of $M^n$:
\[
(1+2S)^{2}\;\geq\;4n\Vert Q\Vert^{2}.
\]
On the other hand, the Cauchy--Schwarz inequality \eqref{e9} gives
$4n\Vert Q\Vert^{2}\geq 4S^{2}$. Combining the two inequalities,
\[
(1+2S)^{2}\geq 4S^{2},
\qquad\text{that is,}\qquad
1+4S\geq0 ,
\]
so that $S\geq-\frac14$ and hence
\[
1+2S\;\geq\;\frac12\;>\;0
\]
everywhere on $M^{n}$. By Vieta's formulas,
\[
r_1r_2=\frac{\|Q\|^2}{n}\ge0,
\qquad
r_1+r_2=\frac{1+2S}{n}>0.
\]
Since the roots are real and have a non-negative product,
they have the same sign (or one of them is zero). Because
their sum is positive, both roots must be non-negative.
As $\lambda$ is one of these roots, we conclude that
\[
\lambda\ge0.
\]
Hence, the soliton is non-expanding; that is, it is either shrinking or steady.

Finally, if $\lambda=0$, then the quadratic gives
$\Vert Q\Vert^{2}=0$, so $Q=0$ and, by \eqref{eq:nablaxi}, $\nabla\xi=0$.
Hence $\xi$ is parallel. Since $\xi$ is a non-vanishing parallel vector
field, the de Rham decomposition theorem implies that
\[
M^n=N^{n-1}\times(a,b),
\]
where $N$ has finite volume. In dimension three, this yields
\[
M^3=\mathbb{R}^3/\Gamma_1 .
\]
\end{proof}

\end{document}
